%% ╔══════════════════════════════════════════════════════════════════╗
%% ║         CORIOTLAB — HOJA DE VIDA — EJEMPLO COMPLETO              ║
%% ║         Datos reales, usados con autorización de su titular      ║
%% ╚══════════════════════════════════════════════════════════════════╝
%%
%% COMPILAR: xelatex -interaction=nonstopmode hoja_de_vida_ejemplo_completa.tex
%%           (doble pasada obligatoria)
%%
%% Este ejemplo NO usa datos ficticios: fue llenado con la hoja de vida
%% real de Santiago Cano Molina (Ingeniero Mecatrónico, ITM), a partir de
%% su CvLAC, Acta de Grado y certificados de cursos/ponencias/prototipos,
%% con su autorización explícita para usarlos como ejemplo de referencia
%% de esta plantilla.
%%
%% Nota sobre la sección "Producción Técnica": este ejemplo no incluye
%% publicaciones tipo artículo/revista porque el titular aún no las tiene
%% registradas; en su lugar se muestra su producción técnica real
%% (prototipos), reutilizando el mismo comando \Publicacion. Si su caso sí
%% tiene publicaciones en revista, simplemente use \Publicacion tal como
%% se documenta en plantilla_hoja_de_vida.tex.

%% ═══════════════════════════════════════════════════════════════════
%% ZONA DE CONFIGURACIÓN — únicos campos a editar
%% ═══════════════════════════════════════════════════════════════════

%% --- Identificación ---
\newcommand{\CVNombre}{Santiago Cano Molina}
\newcommand{\CVTitulo}{Ingeniero Mecatrónico}
\newcommand{\CVCorreo}{santiagocano298366@correo.itm.edu.co}
\newcommand{\CVTelefono}{+57 300 265 4474}
\newcommand{\CVCiudad}{Medellín, Colombia}
\newcommand{\CVLinkedIn}{https://www.linkedin.com/in/santiago-cano-molina}
\newcommand{\CVORCID}{0009-0006-9149-0800}
\newcommand{\CVGitHub}{}
\newcommand{\CVGoogleScholar}{https://scholar.google.com/citations?user=-F0yRQYAAAAJ&hl=es}
\newcommand{\CVResearchGate}{https://www.researchgate.net/profile/Santiago-Cano-Molina-2}
\newcommand{\CVFecha}{Julio de 2026}

%% ═══════════════════════════════════════════════════════════════════

\documentclass[10pt, letterpaper]{article}

%% --- Idioma
\usepackage[spanish, es-nodecimaldot, es-noshorthands]{babel}

%% --- Fuentes (nombres de familia Windows — instalar con setup_coriotlab.ps1)
\usepackage{fontspec}
\setmainfont{Inter}
\newfontfamily\FuenteTitulo{MuseoModerno}
\newfontfamily\FuenteCodigo{Space Mono}

%% --- Geometría — más compacta que los informes: el CV no usa membrete
%%     de página completa, así que se aprovecha más ancho de contenido.
\usepackage[
  top=2.1cm, bottom=2.3cm,
  left=2.2cm, right=2.2cm,
  footskip=1.1cm
]{geometry}

%% --- Colores corporativos CORIOTLAB (idénticos al Libro de Marca)
\usepackage[table]{xcolor}
\definecolor{AzulITM}     {HTML}{102D69}
\definecolor{Magenta}     {HTML}{C14894}
\definecolor{AzulDigital} {HTML}{56ACDE}
\definecolor{GrisPizarra} {HTML}{2F2F2F}
\definecolor{GrisClaro}   {HTML}{F5F5F5}
\definecolor{GrisMedio}   {HTML}{888888}
\definecolor{GrisLinea}   {HTML}{DDDDDD}
\definecolor{Blanco}      {HTML}{FFFFFF}
\definecolor{VerdeTarea}  {HTML}{1A6B3C}

%% --- Gráficos — grffile permite rutas con espacios y puntos múltiples,
%%     necesario porque las carpetas del kit de marca usan ese formato
%%     (ej. "2. Logos/2. Logo Full Color/2. Horizontal/...")
\usepackage{graphicx}
\usepackage{grffile}
\usepackage{tikz}

%% --- Tablas y alineación de entradas
\usepackage{tabularx}
\usepackage{array}
\usepackage{colortbl}
\usepackage{booktabs}

%% --- Listas
\usepackage{enumitem}

%% --- Campos opcionales condicionales (DOI, LinkedIn, ORCID, GitHub, etc.)
\usepackage{etoolbox}

%% --- Secciones y espaciado
\usepackage{titlesec}
\usepackage{needspace}
\usepackage[hidelinks]{hyperref}
\hypersetup{urlcolor=AzulDigital}

%% --- Pie de página discreto (fecha + numeración si el CV excede 1 página)
\usepackage{fancyhdr}

\usepackage{parskip}
\setlength{\parskip}{4pt}
\setlength{\parindent}{0pt}
\renewcommand{\arraystretch}{1.3}

%% ─────────────────────────────────────────────────────────────────
%% PIE DE PÁGINA — regla delgada + fecha (izq) + página (der)
%% ─────────────────────────────────────────────────────────────────
\pagestyle{fancy}
\fancyhf{}
\renewcommand{\headrulewidth}{0pt}
\renewcommand{\footrulewidth}{0.4pt}
\renewcommand{\footrule}{{\color{GrisLinea}\hrule height\footrulewidth}}
\fancyfoot[L]{{\FuenteCodigo\tiny\color{GrisMedio}\CVFecha}}
\fancyfoot[R]{{\small\color{GrisMedio}\thepage}}

%% ─────────────────────────────────────────────────────────────────
%% ESTILO DE SECCIÓN
%% ─────────────────────────────────────────────────────────────────
\newcommand{\SeccionCV}[1]{%
  \needspace{3\baselineskip}%
  \vspace{1.0em}%
  \noindent{\textcolor{AzulITM}{\rule{2.4mm}{2.4mm}}}\hspace{0.5em}%
  {\FuenteTitulo\large\color{AzulITM}\bfseries #1}\\[0.2em]%
  {\color{AzulDigital}\rule{\linewidth}{1.1pt}}%
  \vspace{0.4em}%
}

%% ─────────────────────────────────────────────────────────────────
%% ENCABEZADO DE ENTRADA — fila título (izq) + fechas (der)
%% ─────────────────────────────────────────────────────────────────
\newcommand{\ItemCV}[2]{%
  \par\noindent\begin{tabularx}{\linewidth}{@{}X r@{}}
    {\FuenteTitulo\normalsize\color{GrisPizarra}\bfseries #1} &
    {\FuenteCodigo\small\color{GrisMedio}#2}\\
  \end{tabularx}\par\vspace{0.1em}%
}

\setlist[itemize]{label={\color{AzulITM}\textbf{>}},
  leftmargin=1.6em, labelsep=0.5em, align=left, itemsep=2pt, topsep=3pt}

%% ─────────────────────────────────────────────────────────────────
%% SISTEMA DE ETIQUETAS — Nivel de formación académica (sección 4A)
%% ─────────────────────────────────────────────────────────────────
\newcommand{\NivelFormacion}[1]{%
  \colorbox{AzulITM!12}{\color{AzulITM}\small\bfseries\strut\ #1\ }}
\newcommand{\FormacionEnCurso}{%
  \colorbox{AzulDigital!15}{\color{AzulDigital!70!black}\small\bfseries\strut\ En curso\ }}
\newcommand{\FormacionCulminada}{%
  \colorbox{VerdeTarea!15}{\color{VerdeTarea}\small\bfseries\strut\ Culminado\ }}

%% ─────────────────────────────────────────────────────────────────
%% SISTEMA DE ETIQUETAS — Certificación (sección 4B)
%% ─────────────────────────────────────────────────────────────────
\newcommand{\Certificado}{%
  \colorbox{VerdeTarea!15}{\color{VerdeTarea}\small\bfseries\strut\ Certificado\ }}
\newcommand{\NoCertificado}{%
  \colorbox{GrisMedio!15}{\color{GrisMedio!70!black}\small\bfseries\strut\ Sin certificar\ }}

%% ─────────────────────────────────────────────────────────────────
%% PUBLICACIONES / PRODUCCIÓN TÉCNICA EN FORMATO APA 7 (Opción A)
%% ─────────────────────────────────────────────────────────────────
\newcommand{\Publicacion}[6]{%
  \par\noindent\hangindent=1.2em\hangafter=1
  #1 (#2). #3. \textit{#4}%
  \ifstrempty{#5}{}{, \textit{#5}}%
  .%
  \ifstrempty{#6}{}{\ {\FuenteCodigo\small #6}}%
  \par\vspace{0.5em}%
}

%% ─────────────────────────────────────────────────────────────────
%% CONGRESOS Y PONENCIAS
%% ─────────────────────────────────────────────────────────────────
\newcommand{\Ponencia}[5]{%
  \par\noindent\hangindent=1.2em\hangafter=1
  \textbf{#1}. \textit{#2}. #3, #4 (#5).
  \par\vspace{0.5em}%
}

%% =================================================================
%%  DOCUMENTO
%% =================================================================
\begin{document}
\color{GrisPizarra}

%% ─────────────────────────────────────────────────────────────────
%% ENCABEZADO
%% ─────────────────────────────────────────────────────────────────
\noindent
\begin{minipage}[t]{0.72\linewidth}
  {\FuenteTitulo\fontsize{24}{28}\selectfont\color{AzulITM}\bfseries \CVNombre}\\[0.15em]
  {\large\color{GrisPizarra}\CVTitulo}
\end{minipage}%
\hfill
\begin{minipage}[t]{0.26\linewidth}
  \raggedleft
  \includegraphics[width=0.85\linewidth]%
    {../705_Kit de Marca Coriotlab/2. Logos/2. Logo Full Color/2. Horizontal/Coriotlab_Logo_Hor_FullColor.png}
\end{minipage}

\vspace{0.6em}

%% --- Línea de contacto: correo, teléfono, ciudad ---
\newcommand{\CVSep}{\ \ \textcolor{AzulITM}{\textbf{<}}\ \ }
\noindent{\small\color{GrisMedio}
  \CVCorreo
  \ifdefempty{\CVTelefono}{}{\CVSep\CVTelefono}%
  \CVSep\CVCiudad
}

%% --- Línea de redes/identificadores académicos (separada del contacto) ---
\newbool{CVMostrarRedes}
\ifdefempty{\CVLinkedIn}{}{\booltrue{CVMostrarRedes}}
\ifdefempty{\CVORCID}{}{\booltrue{CVMostrarRedes}}
\ifdefempty{\CVGitHub}{}{\booltrue{CVMostrarRedes}}
\ifdefempty{\CVGoogleScholar}{}{\booltrue{CVMostrarRedes}}
\ifdefempty{\CVResearchGate}{}{\booltrue{CVMostrarRedes}}
\ifbool{CVMostrarRedes}{%
  \vspace{0.25em}

  \noindent{\small\color{GrisMedio}\textit{Redes:}\ \,%
    \def\CVPrimeraRed{1}%
    \ifdefempty{\CVLinkedIn}{}{\href{\CVLinkedIn}{LinkedIn}\def\CVPrimeraRed{0}}%
    \ifdefempty{\CVORCID}{}{\ifnum\CVPrimeraRed=0\CVSep\fi\href{https://orcid.org/\CVORCID}{{\FuenteCodigo ORCID: \CVORCID}}\def\CVPrimeraRed{0}}%
    \ifdefempty{\CVGoogleScholar}{}{\ifnum\CVPrimeraRed=0\CVSep\fi\href{\CVGoogleScholar}{Google Scholar}\def\CVPrimeraRed{0}}%
    \ifdefempty{\CVResearchGate}{}{\ifnum\CVPrimeraRed=0\CVSep\fi\href{\CVResearchGate}{ResearchGate}\def\CVPrimeraRed{0}}%
    \ifdefempty{\CVGitHub}{}{\ifnum\CVPrimeraRed=0\CVSep\fi\href{\CVGitHub}{GitHub}\def\CVPrimeraRed{0}}%
  }%
}{}

\vspace{0.3em}
{\color{GrisLinea}\rule{\linewidth}{0.6pt}}

%% ─────────────────────────────────────────────────────────────────
%% PERFIL PROFESIONAL
%% ─────────────────────────────────────────────────────────────────
\SeccionCV{Perfil Profesional}
Ingeniero mecatrónico egresado del Instituto Tecnológico Metropolitano (ITM),
actualmente cursando la Maestría en Automatización y Control Industrial en la
misma institución. Su trabajo de grado y sus investigaciones posteriores se
han enfocado en el desarrollo de plataformas robóticas móviles e IoT para el
monitoreo y automatización de procesos en el sector agrícola, en terrenos
irregulares. Ha participado como ponente en encuentros científicos nacionales
e internacionales y se desempeña actualmente como apoyo académico
especializado en investigación en la Facultad de Ingenierías del ITM.

%% ─────────────────────────────────────────────────────────────────
%% FORMACIÓN ACADÉMICA
%% ─────────────────────────────────────────────────────────────────
\SeccionCV{Formación Académica}

\ItemCV{Maestría en Automatización y Control Industrial}{Febrero de 2026 -- en curso}
Instituto Tecnológico Metropolitano (ITM)\quad\NivelFormacion{Maestría}\quad\FormacionEnCurso

\vspace{0.7em}

\ItemCV{Ingeniería Mecatrónica}{Febrero de 2020 -- Enero de 2026}
Instituto Tecnológico Metropolitano (ITM)\quad\NivelFormacion{Pregrado}\quad\FormacionCulminada
\begin{itemize}
  \item{} Trabajo de grado: \textit{Desarrollo e implementación de una
        plataforma robótica móvil de bajo costo para aplicaciones en el
        sector agrícola en terrenos irregulares}.
\end{itemize}

%% ─────────────────────────────────────────────────────────────────
%% IDIOMAS
%% ─────────────────────────────────────────────────────────────────
\SeccionCV{Idiomas}
\begin{itemize}
  \item{} \textbf{Español} --- Nativo
  \item{} \textbf{Inglés} --- Habla: Aceptable \textbullet\ Escribe:
        Aceptable \textbullet\ Lee: Bueno \textbullet\ Entiende: Aceptable
\end{itemize}

%% ─────────────────────────────────────────────────────────────────
%% EXPERIENCIA CERTIFICADA
%% ─────────────────────────────────────────────────────────────────
\SeccionCV{Experiencia Certificada}

\ItemCV{Docente cátedra — Extensión Académica}{Abril -- Agosto de 2026}
Instituto Tecnológico Metropolitano (ITM)\quad\Certificado
\begin{itemize}
  \item{} Acompañamiento del proceso formativo del curso de robótica
        (Comuna 12) y del diplomado en inteligencia artificial (Comuna 50),
        en el marco del Proyecto de Inversión de Presupuesto Participativo.
\end{itemize}

\vspace{0.5em}

\ItemCV{Docente cátedra — Extensión Académica}{Marzo -- Julio de 2026}
Instituto Tecnológico Metropolitano (ITM)\quad\Certificado
\begin{itemize}
  \item{} Acompañamiento del proceso formativo en el marco del contrato
        interadministrativo con la Secretaría de Innovación Digital.
\end{itemize}

\vspace{0.5em}

\ItemCV{Apoyo Académico Especializado en Investigación}{Marzo -- Junio de 2026}
Instituto Tecnológico Metropolitano (ITM), Facultad de Ingenierías\quad\Certificado

\vspace{0.5em}

\ItemCV{Prestación de servicios profesionales --- Proyecto SOSAGRO}{Febrero -- Noviembre de 2026}
Instituto Tecnológico Metropolitano (ITM)\quad\Certificado
\begin{itemize}
  \item{} Gestión, acompañamiento y seguimiento administrativo del proyecto
        de regalías Alianza Estratégica SOSAGRO 4.C.
\end{itemize}

\vspace{0.5em}

\ItemCV{Prestación de servicios profesionales --- Detección de Tizón de la Papa}{Diciembre de 2025}
Instituto Tecnológico Metropolitano (ITM)\quad\Certificado
\begin{itemize}
  \item{} Apoyo técnico especializado en el desarrollo, implementación y
        optimización de algoritmos de procesamiento y análisis de datos
        espectrales en infraestructura de cómputo de alto rendimiento (HPC),
        orientados a la mejora de modelos de detección del tizón de la papa.
\end{itemize}

\vspace{0.5em}

\ItemCV{Formador}{Mayo -- Noviembre de 2025}
Convenio interadministrativo Comfenalco Antioquia -- ITM\quad\Certificado
\begin{itemize}
  \item{} Formación en robótica en el marco del convenio interadministrativo
        entre Comfenalco Antioquia y el Instituto Tecnológico Metropolitano
        (ITM).
\end{itemize}

%% ─────────────────────────────────────────────────────────────────
%% CERTIFICACIONES
%% ─────────────────────────────────────────────────────────────────
\SeccionCV{Certificaciones}

\ItemCV{The 3D Printing Revolution}{Noviembre de 2025}
University of Illinois Urbana-Champaign / Coursera\quad{\FuenteCodigo\small\color{GrisMedio}coursera.org/verify/3E8C7CVAHEBQ}

\vspace{0.5em}

\ItemCV{Fundamentos Básicos y Aplicaciones Específicas de la IA Generativa}{Noviembre de 2025}
Beca SER ANDI -- Nodo EAFIT

\vspace{0.5em}

\ItemCV{Bootcamp de Análisis de Datos, Nivel Intermedio (159 horas)}{Octubre de 2025}
Talento Tech (MinTIC) -- IU Training: U. de Antioquia, U. de Caldas, Ubicua Technology

\vspace{0.5em}

\ItemCV{Fundamentos de Machine Learning}{Octubre de 2025}
Platzi\quad{\FuenteCodigo\small\color{GrisMedio}Cód. 536c6e18-b531-4c21-9835-287cfbf9d0af}

\vspace{0.5em}

\ItemCV{Álgebra Lineal Aplicada para Machine Learning}{Octubre de 2025}
Platzi\quad{\FuenteCodigo\small\color{GrisMedio}Cód. a4721bc3-e758-4c02-9363-9c88528d002a}

\vspace{0.5em}

\ItemCV{Inmersión IA (con Google Gemini)}{Septiembre de 2025}
Alura

\vspace{0.5em}

\ItemCV{Preparación Certificación AWS AI Practitioner}{Agosto de 2025}
Alcaldía de Medellín -- Talento Tech

\vspace{0.5em}

\ItemCV{Técnicas de análisis de datos para un futuro sostenible en agricultura y medio ambiente}{Octubre de 2024}
Universidad Surcolombiana, II Congreso Internacional USRA e Ingeniería Agrícola

\vspace{0.5em}

\ItemCV{Introduction to Bionic Prostheses, 3D Biomechanical Prostheses y Bionic Exoskeletons (3 webinars)}{Septiembre de 2024}
P4H Bionics Online Academy

\vspace{0.5em}

\ItemCV{Python para no matemáticos: de 0 hasta reconocimiento facial}{Marzo de 2024}
Udemy\quad{\FuenteCodigo\small\color{GrisMedio}Cód. UC-7735021d-de59-47f1-8af5-59c31652f210}

\vspace{0.5em}

\ItemCV{Robótica: Programación de Robots con ROS}{Agosto de 2023}
Udemy\quad{\FuenteCodigo\small\color{GrisMedio}Cód. UC-257a85f3-cad3-4ee1-9d03-dbc9deed6187}

\vspace{0.5em}

\ItemCV{MATLAB Fundamentals}{Noviembre de 2022}
MathWorks Training Services

\vspace{0.5em}

\ItemCV{Introducción a Python con énfasis en manejo de datos e IA}{Agosto -- Septiembre de 2022}
Instituto Tecnológico Metropolitano (ITM), Extensión Académica

\vspace{0.5em}

\ItemCV{Programa Tecnoacademia (13 cursos y talleres, 702 horas acumuladas)}{2014 -- 2021}
Servicio Nacional de Aprendizaje (SENA), Regional Antioquia
\begin{itemize}
  \item{} Introducción a la Programación: Algoritmos, Diseño y Prototipado
        (2 ediciones), Dibujo Técnico, Ciencias Básicas en Matemáticas y
        Física, Herramientas Metodológicas en Investigación Aplicada (CTeI),
        Aplicación de Buenas Prácticas y Manejo de Reactivos en Laboratorio
        Químico, Ilustración con Adobe Illustrator, Metodología de Marco
        Lógico para Proyectos Públicos, Manejo de Microsoft Office 2016
        (Excel), Habilidades para la Vida y Sensibilización Tecnoacademia.
\end{itemize}

%% ─────────────────────────────────────────────────────────────────
%% PRODUCCIÓN TÉCNICA
%% Nota: el titular aún no tiene publicaciones tipo artículo/revista
%% registradas; se muestra en su lugar la producción técnica real
%% (registros de software y prototipos), reutilizando el comando
%% \Publicacion.
%% ─────────────────────────────────────────────────────────────────
\SeccionCV{Producción Técnica}

\Publicacion{Cano Molina, S., y Hernández Palacio, L. F.}{2025}%
  {CULTIVA: sistema inteligente de recomendación agrícola}%
  {Registro de software --- Dirección Nacional de Derecho de Autor (DNDA), Colombia}%
  {}{Reg. 13-108-391}

\Publicacion{Cano Molina, S., y Hernández Palacio, L. F.}{2025}%
  {RiegoSmart: control Bluetooth, posicionamiento mediante IMU y
  publicación a Ubidots}%
  {Registro de software --- Dirección Nacional de Derecho de Autor (DNDA), Colombia}%
  {}{Reg. 13-109-62}

\Publicacion{Cano Molina, S., Hernández Palacio, L. F., y Mejía Herrera, M.}{2025}%
  {Sistema de captación de agua atmosférica}%
  {Prototipo industrial --- Instituto Tecnológico Metropolitano (ITM), Colombia}{}{}

\Publicacion{Cano Molina, S., Hernández Palacio, L. F., y Mejía Herrera, M.}{2025}%
  {Sistema portátil de captación de agua atmosférica}%
  {Prototipo industrial --- Instituto Tecnológico Metropolitano (ITM), Colombia}{}{}

%% ─────────────────────────────────────────────────────────────────
%% CONGRESOS Y PONENCIAS
%% ─────────────────────────────────────────────────────────────────
\SeccionCV{Congresos y Ponencias}

\Ponencia{Desarrollo e implementación de una plataforma robótica móvil de
  bajo costo para aplicaciones en el sector agrícola en terrenos irregulares}%
  {11.\textsuperscript{a} Semana Internacional de Ciencia, Tecnología e
  Innovación --- Universidad Francisco de Paula Santander}%
  {Ponente}{San José de Cúcuta, Colombia}{Octubre de 2024}

\Ponencia{Desarrollo de una plataforma robótica móvil IoT para el monitoreo
  de variables ambientales en entornos agrícolas}%
  {Seminario Científico, Reto Autoinnova 2024 --- Universidad Central
  ``Marta Abreu'' de Las Villas (UCLV)}%
  {Ponente}{Santa Clara, Cuba}{Abril de 2025}

\Ponencia{Sistema inteligente de control de bombas para riego en torres de
  cultivo}%
  {XPOILERS: Festival de Stand Up Científicos --- Iniciación Científica /
  Instituto Tecnológico Metropolitano (ITM)}%
  {Ponente}{Bogotá D.\,C., Colombia}{Abril -- Mayo de 2026}

\Ponencia{Desarrollo de sistema de navegación autónoma con tecnología de
  banda ultra ancha}%
  {II Congreso Internacional USRA e Ingeniería Agrícola --- Universidad
  Surcolombiana}%
  {Ponente}{Neiva, Colombia}{Octubre de 2024}
\noindent{\small\color{GrisMedio}Coautores: Lucas Felipe Hernández Palacio,
  Mateo Mejía Herrera, Juan Sebastián Botero Valencia.}
\vspace{0.5em}

\Ponencia{Participación con proyecto de semillero de investigación}%
  {XX Encuentro Departamental de Semilleros de Investigación, Nodo
  Antioquia --- RedCOLSI}%
  {Ponente}{Virtual}{Junio de 2021}

%% ─────────────────────────────────────────────────────────────────
%% RECONOCIMIENTOS
%% ─────────────────────────────────────────────────────────────────
\SeccionCV{Reconocimientos}

\ItemCV{Primer lugar --- Reto Autoinnova 2024}{Diciembre de 2024}
Facultad de Ingenierías, Instituto Tecnológico Metropolitano (ITM)
\begin{itemize}
  \item{} Proyecto \textit{Desarrollo de una plataforma robótica móvil IoT
        para el monitoreo de variables ambientales en invernaderos
        agrícolas}, junto con Lucas Felipe Hernández Palacio.
\end{itemize}

\vspace{0.5em}

\ItemCV{Reconocimiento --- ROBORAVE México, 3.\textsuperscript{a} edición}{Septiembre de 2016}
ROBORAVE International

\vspace{0.5em}

\ItemCV{Evaluador, XIV y XV Muestra Institucional de Proyectos de Investigación}{2024 -- 2025}
Institución Educativa Colegio Loyola para la Ciencia y la Innovación

\end{document}
